\documentclass{article}
\usepackage{amsfonts,amsthm,amsmath,amssymb,mathtools}
\usepackage{mathrsfs}
\usepackage{enumitem}
\usepackage{tikz-cd}

% Chapter
\newcounter{chapter}
\setcounter{chapter}{1}

% Section
\setcounter{section}{4}

\newtheorem{thm}{Theorem}
\newenvironment{theorem}[1]{
	\renewcommand{\thethm}{#1}
	\begin{thm}
		}{\end{thm}}

\newtheorem{lma}{Lemma}
\newenvironment{lemma}[1]{
	\renewcommand{\thelma}{#1}
	\begin{lma}
		}{\end{lma}}

\theoremstyle{definition}
\newtheorem{ex}{}[section]
\newenvironment{exercise}[1]{
  \renewcommand{\theex}{\arabic{section}.\arabic{ex}}
	\begin{ex}
		}{\end{ex}}

\title{Exercises {\Roman{chapter}}.{\arabic{section}}}
\author{Connor Mooneyhan}
\date{June 15, 2024}

\begin{document}

\maketitle

\begin{ex}
	Composition is defined for \textit{two} morphisms.
	If more than two morphisms are given, e.g., \begin{equation*}
		\begin{tikzcd}
			A \arrow[r, "f"] & B \arrow[r, "g"] & C \arrow[r, "h"] & D \arrow[r, "i"] & E,
		\end{tikzcd}
	\end{equation*}
	then one may compose them in several ways, for example: \begin{equation*}
		(ih)(gf), \hspace{1em} (i(hg))f, \hspace{1em} i((hg)f), \hspace{1em} \mathrm{etc.}
	\end{equation*}
	so that at every step one is only composing two morphisms.
	Prove that the result of any such nested composition is independent of the placement of the parentheses.
	(Hint: Use induction on $n$ to show that any such choice for $f_nf_{n-1}\cdots f_1$ equals \begin{equation*}
		((\cdots((f_nf_{n-1})f_{n-2})\cdots)f_1).
	\end{equation*}
	Carefully working out the case $n = 5$ is helpful.)
\end{ex}
\begin{proof}
	Note that the cases of $n = 0$, $n = 1$, $n = 2$ are trivially true (parentheses can only be applied in at most 1 way!), so we observe that $n = 3$ is simply an expression of the fact that morphisms are associative.
	This is true in any category, and it serves as our "true" base case.
	As a matter of notation, we will denote the cumbersome \begin{equation*}
		((\cdots((f_nf_{n-1})f_{n-2})\cdots)f_1)
	\end{equation*}
	as simply \begin{equation*}
		f_n f_{n-1} \cdots f_1,
	\end{equation*}
	dropping all parentheses except where they are either required or clarifying.
	In this case, composition is intended to be read "left-to-right."

	Let $n > 3$ such that the theorem holds for all $m \leq n - 1$.
	Then given any valid arrangement of parentheses, since at every step one is only composing two morphisms, there must exist (proven in the next paragraph) at least one pair $f_kf_{k - 1}$ where $n \geq k > 1$ such that when all parentheses are written out, $(f_kf_{k-1})$ is a part of the product.
	We will call such a pair a \textit{primitive pair}.
	For example, in the product \begin{equation*}
		(f_5(f_4((f_3f_2)f_1))),
	\end{equation*}
	$(f_3f_2)$ is a primitive pair.

	We now provide a brief proof of its existence.
	It is clear that such a pair exists for $n = 3$ since the only two arrangements of parentheses are $((f_3f_2)f_1)$ and $(f_3(f_2f_1))$, both of which feature a primitive pair.
	If $n > 3$, then in the case that $(f_nf_{n-1})$ is a primitive pair, we are done, so assume that $f_n$ is not part of a primitive pair.
	Then $f_n$ is being multiplied by some parentheses arrangement of the product $f_{n-1}f_{n-2}\cdots f_1$, which itself contains a primitive pair by the induction hypotheses.

	We conclude, then, that each arrangement of parentheses as specified does indeed contain at least one primitive pair, and since $n$ is finite, there must be a "first" or "leftmost" primitive pair.
	Fix $k$ so that this leftmost primitive pair is $(f_kf_{k-1})$.
	Then we can consider $(f_k f_{k - 1})$ as one morphism, making our product a parentheses arrangement of \begin{equation*}
		f_n f_{n-1} \cdots (f_kf_{k - 1}) \cdots f_1,
	\end{equation*}
	which is a product of $n - 1$ elements, and thus by our induction hypothesis is equal to \textit{exactly} \begin{equation*}
		f_n f_{n-1} \cdots (f_kf_{k - 1}) \cdots f_1.
	\end{equation*}
	We can now transform this product, adding parentheses when useful: \begin{align*}
		  & f_n f_{n-1} \cdots (f_kf_{k - 1}) \cdots f_1                \\
		= & (f_n f_{n-1} \cdots (f_kf_{k - 1})) \cdots f_1              \\
		= & (f_n f_{n-1} \cdots (f_kf_{k - 1})) (f_{k - 2} \cdots f_1),
	\end{align*}
	where the second equality comes from the induction hypothesis.
	When $k > 2$, we can use the hypothesis again to simplify this to \begin{align*}
		  & (f_n f_{n-1} \cdots f_kf_{k - 1}) (f_{k - 2} \cdots f_1)         \\
		= & (\cdots((f_n f_{n-1} \cdots f_kf_{k - 1}) f_{k - 2}) \cdots f_1) \\
		= & f_n f_{n-1} \cdots f_kf_{k - 1} f_{k - 2} \cdots f_1,
	\end{align*}
	which is our desired result.
	If $k = 2$, then we use a similar trick: \begin{align*}
		  & f_n f_{n-1} \cdots (f_2 f_1)       \\
		= & (f_n f_{n-1} \cdots f_3) (f_2 f_1) \\
		= & (((f_n f_{n-1}\cdots f_3)f_2)f_1)  \\
		= & f_n f_{n-1}\cdots f_3f_2f_1.
	\end{align*}
\end{proof}

\begin{ex}
	In Example 3.3 we have seen how to construct a category from a set endowed with a relation, provided this latter is reflexive and transitive.
	For what types of relations is the corresponding category a groupoid?
\end{ex}
\begin{proof}
	Reflexivity and transitivity are still needed to guarantee the existence of identities and composition respectively, but we only need to add one axiom to form a groupoid: symmetry.
	Thus, our relation would be an equivalence relation.
	We now proceed to prove that a category formed by a reflexive and transitive relation is a groupoid if and only if it is symmetric.
	Let $\mathsf{C}$ be our category and $\sim$ our relation.

	($\implies$) Given $A$ and $B$ two objects of $\mathsf{C}$ such that $A \sim B$, by definition $\exists f \in \mathrm{Hom}_{\mathsf{C}}(A, B)$.
	We also know, since $\mathsf{C}$ is a groupoid, that $f$ is an isomorphism and thus has an inverse, so $\exists f^{-1} \in \mathrm{Hom}_{\mathsf{C}}(B, A)$, meaning that $B \sim A$.
	Thus $\sim$ is symmetric.

	($\impliedby$) Given $f \in \mathrm{Hom}_{\mathsf{C}}(A, B)$ for some objects $A$ and $B$, we know by definition that $A \sim B$.
	Since $\sim$ is symmetric, we also have $B \sim A$, and thus $\exists g \in \mathrm{Hom}_{\mathsf{C}}(B, A)$.
	Thus $gf \in \mathrm{End}_{\mathsf{C}}(A)$, but $\mathrm{End}_{\mathsf{C}}(A)$ has only one element, namely $1_A$, so $gf = 1_A$.
	A similar proof shows that $fg = 1_B$, so we conclude that $g$ is the inverse of $f$, and thus $f$ is an isomorphism.
	Since the choice of $f$ was arbitrary, $\mathsf{C}$ is a groupoid.
\end{proof}

\begin{ex}
	Let $A, B$ be objects of a category $\mathsf{C}$, and let $f \in \mathrm{Hom}_{\mathsf{C}}(A, B)$ be a morphism.
	\begin{itemize}
		\item Prove that if $f$ has a right-inverse, then $f$ is an epimorphism.
		\item Show that the converse does not hold, by giving an explicit example of a category and an epimorphism without a right-inverse.
	\end{itemize}
\end{ex}
\begin{proof}
	First, assume $f$ has a right-inverse $g \in \mathrm{Hom}_{\mathsf{C}}(B, A)$.
	Given $f_1, f_2 \in \mathrm{Hom}_{\mathsf{C}}(B, Z)$ with $Z$ some object of $\mathsf{C}$ such that $f_1 f = f_2 f$, then \begin{equation*}
		f_1 f g = f_2 f g \implies f_1 (f g) f_2 (f g) \implies f_1 1_B = f_2 1_B \implies f_1 = f_2.
	\end{equation*}

	On the other hand, take the example of the category generated (as in Example 3.3) by the relation $\leq$ on $\mathbb{N}$.
	There is a morphism $f \in \mathrm{Hom}(2, 3)$, and as pointed out in the present section, it is an epimorphism just like every other morphism in this category.
	However, $\mathrm{Hom}(3, 2) = \emptyset$, so there cannot be a right-inverse.
\end{proof}

\begin{ex}
	Prove that the composition of two monomorphisms is a monomorphism.
	Deduce that one can define a subcategory $\mathsf{C}_{\mathrm{mono}}$ of a category $\mathsf{C}$ by taking the same objects as in $\mathsf{C}$ and defining $\mathrm{Hom}_{\mathsf{C}_{\mathrm{mono}}}(A, B)$ to be the subset of $\mathrm{Hom}_{\mathsf{C}}(A, B)$ consisting of monomorphisms, for all objects $A, B$.
	(Cf. Exercise 3.8; of course, in general $\mathsf{C}_{\mathrm{mono}}$ is not full in $\mathsf{C}$.)
	Do the same for epimorphisms.
	Can you define a subcategory $\mathsf{C}_{\mathrm{nonmono}}$ of $\mathsf{C}$ by restricting to morphisms that are \textit{not} monomorphisms?
\end{ex}
\begin{proof}
  Given two monomorphisms $f \in \mathrm{Hom}_{\mathsf{C}}(A, B)$ and $g \in \mathrm{Hom}_{\mathsf{C}}(B, C)$ for objects $A, B, C$ in a category $\mathsf{C}$, we can deduce easily that given $h_1, h_2 \in \mathrm{Hom}_{\mathsf{C}}(Z, A)$ for some object $Z$, \begin{align*}
		         & (gf)h_1 = (gf)h_2 \\
		\implies & g(fh_1) = g(fh_2) \\
		\implies & fh_1 = fh_2       \\
		\implies & h_1 = h_2.
	\end{align*}
  Thus the composition of monomorphisms is indeed a monomorphism. 
  Also, identities are monomorphisms, so $\mathsf{C}_{\mathrm{mono}}$ is a category.

  On the other hand, since identities are monomorphisms, they are not objects in the proposed $\mathsf{C}_{\mathrm{nonmono}}$, so $\mathsf{C}_{\mathrm{nonmono}}$ cannot form a subcategory.
\end{proof}

\begin{ex}
  Give a concrete description of monomorphisms and epimorphisms in the category $\mathsf{MSet}$ you constructed in Exercise 3.9.
  (Your answer will depend on the notion of morphism you defined in that exercise!)
\end{ex}
\begin{proof}
  Monomorphisms here are still monomorphisms of sets, just sets of equivalence classes.
  If you want to interpret the monomorphisms in terms of "the original set," then you can view it as a function being a monomorphism if it sends \textit{non-related} (by $\sim$) elements to different elements.
  For epimorphisms, the concept stays more or less the same, where a function is an epimorphism if every element of the codomain is sent to by some equivalence class of related elements.
\end{proof}

\end{document}
